% Options for packages loaded elsewhere
\PassOptionsToPackage{unicode}{hyperref}
\PassOptionsToPackage{hyphens}{url}
\PassOptionsToPackage{dvipsnames,svgnames*,x11names*}{xcolor}
%
\documentclass[
  english,
  a4paper,
  DIV=12,
  footlines=2.1,
  usegeometry=true]{scrartcl}
\usepackage{lmodern}
\usepackage{setspace}
\setstretch{1.1}
\usepackage{amssymb,amsmath}
\usepackage{ifxetex,ifluatex}
\ifnum 0\ifxetex 1\fi\ifluatex 1\fi=0 % if pdftex
  \usepackage[T1]{fontenc}
  \usepackage[utf8]{inputenc}
  \usepackage{textcomp} % provide euro and other symbols
\else % if luatex or xetex
  \usepackage{unicode-math}
  \defaultfontfeatures{Scale=MatchLowercase}
  \defaultfontfeatures[\rmfamily]{Ligatures=TeX,Scale=1}
  \setmainfont[]{Charis SIL}
  \setsansfont[]{Charis SIL Compact}
  \ifxetex
    \usepackage[space]{xeCJK}
    \setCJKmainfont[]{Noto Serif CJK TC}
  \fi
  \ifluatex
    \usepackage[]{luatexja-fontspec}
    \setmainjfont[]{Noto Serif CJK TC}
  \fi
\fi
% Use upquote if available, for straight quotes in verbatim environments
\IfFileExists{upquote.sty}{\usepackage{upquote}}{}
\IfFileExists{microtype.sty}{% use microtype if available
  \usepackage[]{microtype}
  \UseMicrotypeSet[protrusion]{basicmath} % disable protrusion for tt fonts
}{}
\usepackage{xcolor}
\IfFileExists{xurl.sty}{\usepackage{xurl}}{} % add URL line breaks if available
\IfFileExists{bookmark.sty}{\usepackage{bookmark}}{\usepackage{hyperref}}
\hypersetup{
  pdftitle={Horrordatei zum Testen vom neuen Textsatz},
  pdflang={en},
  pdfkeywords={Entangled Religions, religious contact, Markdown, OES},
  colorlinks=true,
  linkcolor=NavyBlue,
  filecolor=Maroon,
  citecolor=Blue,
  urlcolor=Blue,
  pdfcreator={LaTeX via pandoc}}
\urlstyle{same} % disable monospaced font for URLs
\usepackage{longtable,booktabs}
% Correct order of tables after \paragraph or \subparagraph
\usepackage{etoolbox}
\makeatletter
\patchcmd\longtable{\par}{\if@noskipsec\mbox{}\fi\par}{}{}
\makeatother
% Allow footnotes in longtable head/foot
\IfFileExists{footnotehyper.sty}{\usepackage{footnotehyper}}{\usepackage{footnote}}
\makesavenoteenv{longtable}
\usepackage{graphicx}
\makeatletter
\def\maxwidth{\ifdim\Gin@nat@width>\linewidth\linewidth\else\Gin@nat@width\fi}
\def\maxheight{\ifdim\Gin@nat@height>\textheight\textheight\else\Gin@nat@height\fi}
\makeatother
% Scale images if necessary, so that they will not overflow the page
% margins by default, and it is still possible to overwrite the defaults
% using explicit options in \includegraphics[width, height, ...]{}
\setkeys{Gin}{width=\maxwidth,height=\maxheight,keepaspectratio}
\usepackage[normalem]{ulem}
% Avoid problems with \sout in headers with hyperref
\pdfstringdefDisableCommands{\renewcommand{\sout}{}}
\setlength{\emergencystretch}{3em} % prevent overfull lines
\providecommand{\tightlist}{%
  \setlength{\itemsep}{0pt}\setlength{\parskip}{0pt}}
\setcounter{secnumdepth}{-2}

% Page style, header and footer
\usepackage[headsepline]{scrlayer-scrpage}
\pagestyle{scrheadings}
\lohead{\textsc{Reiker -- Horrordatei zum Testen vom neuen Textsatz}}
\rohead{\textit{Entangled Religions} 9 (2019)}
%\cofoot[]{\thepage}
\cofoot[]{}

%% ER Settings
%% Page margin

% We don’t use the page foot except on the first page.
% This leads to a large bottom margin on most pages.
% To solve this, we’re setting the bottom margin manually
% and set footskip to 0pt. This completely merges the
% foot into the text area.
% We’re then using \enlargethispage{-5\baselineskip} to
% create space for the foot on the first page.
\usepackage[bottom=2.8cm,footskip=0pt]{geometry}

%% Force loading of required packages
\usepackage{booktabs,graphicx}

%% Bookmarks/Outline
\usepackage{bookmark}

%% Text settings
\tolerance=1000
\widowpenalty=10000
\clubpenalty=10000

%% Colors
\definecolor{parnum}{gray}{0.5}

%% Figure captions
\KOMAoptions{captions=tableabove}
\addtokomafont{caption}{\normalfont\small}
\addtokomafont{captionlabel}{\bfseries}
\renewcommand*{\captionformat}{\quad}

%% page headings
\setkomafont{pageheadfoot}{%
\normalfont\normalcolor\small}
\setkomafont{author}{\large\scshape}
\setkomafont{date}{\large}

%% languages/scripts
\newfontfamily\hebrewfont[Script=Hebrew]{FreeSerif}
\newfontfamily\arabicfont[Script=Arabic]{Scheherazade}
\newfontfamily\syriacfont[Script=Syriac]{Noto Sans Syriac Estrangela}
\newfontfamily\greekfont[Script=Greek]{Gentium Plus}
\newfontfamily\copticfont[Script=Coptic]{Noto Sans Coptic}
\newfontfamily\ethiopicfont[Script=Ethiopic]{Noto Serif Ethiopic}
\newfontfamily\amharicfont[Script=Ethiopic]{Noto Serif Ethiopic}

%% enumerations

\renewcommand{\labelenumi}{\arabic{enumi}.}
\renewcommand{\labelenumii}{\arabic{enumi}.\arabic{enumii}}
\renewcommand{\labelenumiii}{\arabic{enumi}.\arabic{enumii}.\arabic{enumiii}}

%% figure positioning
\makeatletter
\def\fps@figure{tbp}
\makeatother

%% set floats above footnotes
\usepackage[bottom]{footmisc}

%% End ER Settings

\ifxetex
  % Load polyglossia as late as possible: uses bidi with RTL langages (e.g. Hebrew, Arabic)
  \usepackage{polyglossia}
  \setmainlanguage[]{english}
  \setotherlanguage[]{hebrew}
  \setotherlanguage[]{english}
  \setotherlanguage[]{syriac}
  \setotherlanguage[]{arabic}
\else
  \usepackage[shorthands=off,hebrew,english,syriac,arabic,main=english]{babel}
  \newcommand{\texthebrew}[2][]{\foreignlanguage{hebrew}{#2}}
  \newenvironment{hebrew}[2][]{\begin{otherlanguage}{hebrew}}{\end{otherlanguage}}
  \newcommand{\textenglish}[2][]{\foreignlanguage{english}{#2}}
  \newenvironment{english}[2][]{\begin{otherlanguage}{english}}{\end{otherlanguage}}
  \newcommand{\textsyriac}[2][]{\foreignlanguage{syriac}{#2}}
  \newenvironment{syriac}[2][]{\begin{otherlanguage}{syriac}}{\end{otherlanguage}}
  \newcommand{\textarabic}[2][]{\foreignlanguage{arabic}{#2}}
  \newenvironment{arabic}[2][]{\begin{otherlanguage}{arabic}}{\end{otherlanguage}}
\fi

%% Footnotes (load after polyglossia!)
\deffootnote{\leftmargin}{1em}{%
\makebox[\leftmargin][l]{\thefootnotemark}%
}

\newlength{\cslhangindent}
\setlength{\cslhangindent}{\leftmargin}
\newenvironment{cslreferences}%
  {\setlength{\parindent}{0pt}%
  \everypar{\setlength{\hangindent}{\cslhangindent}}\ignorespaces}%
  {\par}

\title{Horrordatei zum Testen vom neuen Textsatz}
\usepackage{etoolbox}
\makeatletter
\providecommand{\subtitle}[1]{% add subtitle to \maketitle
  \apptocmd{\@title}{\par {\large #1 \par}}{}{}
}
\makeatother
\subtitle{Untertitel zur Horrordatei}
\author{Julia Reiker\\Centrum für Religionswissenschaftliche Studien,
Ruhr-Universität Bochum}
\date{}

\begin{document}
%% No additional space between sentences
\frenchspacing

\cleardoubleoddpage
\thispagestyle{empty}
\begin{singlespace}

%% Logo
\hfill\includegraphics[height=3\baselineskip]{Logo_ER.pdf}
\noindent\rule{\textwidth}{\heavyrulewidth}

%% Head
\vspace*{.2\baselineskip}
\noindent\begin{minipage}[c]{.13\textwidth}
\scriptsize
\parbox{.2\textwidth}{\includegraphics[height=2\baselineskip,width=\csname Gin@nat@width\endcsname,keepaspectratio]{cc-by.pdf}}
\end{minipage}
\begin{minipage}[c]{.4\textwidth}
\scriptsize
\raggedright
DOI:~\href{https://doi.org/http://doi.org./10.13154/volume.year.author\_shorttitle}{http://doi.org./10.13154/volume.year.author\_shorttitle}\\
License:~\href{http://creativecommons.org/licenses/by/4.0/}{Creative
Commons Attribution 4.0}
\end{minipage}\hfill
\begin{minipage}[c]{.4\textwidth}
\scriptsize
\raggedleft
\textit{Entangled Religions} 9 (2019)\\
\href{https://er.ceres.rub.de/}{er.ceres.rub.de}
\end{minipage}

%% Title
\vspace*{2\baselineskip}
{\usekomafont{title}\LARGE\noindent\raggedright Horrordatei zum Testen
vom neuen Textsatz\par}
\vspace*{.5\baselineskip}
{\usekomafont{subtitle}\Large\noindent\raggedright Untertitel zur
Horrordatei\par}
\vspace{2\baselineskip}
{\usekomafont{author}\noindent Julia Reiker\par}
{\usekomafont{author}\normalfont\itshape\normalsize\noindent Centrum für
Religionswissenschaftliche Studien, Ruhr-Universität Bochum\par}
\vspace{\baselineskip}
\vspace{2\baselineskip}
\end{singlespace}
\begin{quote}
\small
\begin{tabular}{@{}p{\linewidth}@{}}
\toprule
{\bfseries\scshape abstract}\hspace{1em}Lorem ipsum dolor sit amet,
consetetur sadipscing elitr, sed diam nonumy eirmod tempor invidunt ut
labore et dolore magna aliquyam erat, sed diam voluptua. At vero eos et
accusam et justo duo dolores et ea rebum. Stet clita kasd gubergren, no
sea takimata sanctus est Lorem ipsum dolor sit amet. Lorem ipsum dolor
sit amet, consetetur sadipscing elitr, sed diam nonumy eirmod tempor
invidunt ut labore et dolore magna aliquyam erat, sed diam voluptua. At
vero eos et accusam et justo duo dolores et ea rebum. Stet clita kasd
gubergren, no sea takimata sanctus est Lorem ipsum dolor sit amet.\\
\midrule[\heavyrulewidth]
{\bfseries\scshape keywords}\hspace{1em}Entangled Religions, religious
contact, Markdown, OES\\
\bottomrule
\end{tabular}
\end{quote}
\vspace{1.8181\baselineskip} % Since we altered line spacing, we need to adjust the vspace here.

\setstretch{1.1}
\hypertarget{beispiel-fuxfcr-andere-sprachen}{%
\section{Beispiel für andere
Sprachen}\label{beispiel-fuxfcr-andere-sprachen}}

\begin{hebrew}

\leavevmode\marginpar{\textcolor{parnum}{{[}1{]}}}\textenglish{{[}I{]}}(11)
\ldots{} ותען חנה ות{[}אמר{]}. (12) וגואב דאד חנה וגופת נא סייד זני סכת
ג͘אן הום מן ונביד ומוסכיר (13) נא הינגיסתום והמי ריזום תהלי גאן מרא פא
פישי י'י. ותען חנה. תא (14) תא גואב {[}דאדי{]}ש חנה צון ש{[}ניד{]} אין
רא וגופתיש לא אדני. בֿ ציז (15) גופתיש. יכי גופת נא נא אי סייד נא נא הום
מסת. דויום המי כואהום (16) שאדי כי {[}היסת?{]} ישראלן רא וליכין אשה קשת
רוח אנכי. סכת (17) תהליהא וגמי {[}היסת{]} פא דילי מן וכוד יין ˺ושכר ק˹
לא שתיתי. וצון חנה (18) גופת כוישתן רא קשת רוח אנכי. המצונין גלותיאן רא
גופת כי (19) כאשה עזובה ועצובת רוח קראך י'י ואשת נעורים. כי תמאס אמר
(20) אלהיך. וצון כי חנה רא שאד כרד י'י המצונין שאד כונד י'י אין (21)
גלותיאן רא: \textsuperscript{צ͘ה͘ג͘~ובני~ציון~גילו~ושמחו~וג}͘ ואשפוך את
נפשי לפני י'י. גופת כי המי ריזום תהלי (22) גאן מן רא פישי י'י תא מוראדי
מן בידהד. המצונין אין גלות{[}י{]}אן (23) המי גוינד אשפוך לפניו שיחי צרתי
לפניו וג: אלה אזכרה ואש{[}פכ{]}ה (24) עלי נפשי כי אעבור וג:

\end{hebrew}

\leavevmode\marginpar{\textcolor{parnum}{{[}2{]}}}\texthebrew{אלי נאיר
צהרי ענדי דינ' ואחד} (\emph{ʾly nʾyr ṣhri ʿndi dyn' wʾḥd}), ``I owe my
brother-in-law Nāyar one dinar'' (TS 20.137, line 2.6). The term
\emph{ṣihr} refers to in-law in general, and may be either
father-in-law, brother-in-law or son-in-law. In the case of Aśu, an
in-law through matrimonial alliance to her could only be her brother,
for \emph{nāyar}s follow a matrilineal system, where the brother is of
more significance to the woman's kinship relations than the father
(Gough 1961, 352-4).

\leavevmode\marginpar{\textcolor{parnum}{{[}3{]}}}This is concurrent
with the majority of Peshitta manuscripts to this verse:
``\textsyriac{ܡܫܚܐ ܕܪ̈ܫܝܥܐ ܠܪܝܫܝ ܠܐ ܢܕܗܢ}'', ``\emph{mšḥa d'ryš}‛\emph{a
lryšy la ndhn}'' (``oil of the wicked will not anoint my head''), and is
in line with the Septuagint's reading of the verse, too (Carbajosa 2008,
263).

\begin{Arabic}

\leavevmode\marginpar{\textcolor{parnum}{{[}4{]}}}(3) \ldots{} ثم قالوا
\textbf{كي ايل ديعو{[}ث{]}} (4) \textbf{יי'} وفيه قولان الواحد \ldots{}
والاخر ان الله عارف الخفيات عن الناس فمن تعاطا الكلام (6) عليها فقد تكلم
بما ينكره الله عليه وفي مثل ذلك قال داود عليه السلم (7) \textbf{ولوا
هلاختي بغذولوث وبنيفلاوث مماني}

\end{Arabic}

\leavevmode\marginpar{\textcolor{parnum}{{[}5{]}}}This manuscript is
most likely one of the oldest Karaite manuscripts written in Arabic
characters. For its description, see Margoliouth 1899, 207-208; Hoerning
1889, 45-60. For studies concerning Karaite manuscripts of the Bible in
Arabic characters, see Hoerning 1889, Khan 1990. For the importance of
the manuscripts of Yefet b. ʿEli's commentaries written in Arabic
characters, see Ben-Shammai 1976.

\leavevmode\marginpar{\textcolor{parnum}{{[}6{]}}}Beispiel für eine
Aufzählung:

\begin{enumerate}
\item
  \leavevmode\marginpar{\textcolor{parnum}{{[}7{]}}}Opening /
  Introduction (§1--2)
\item
  \leavevmode\marginpar{\textcolor{parnum}{{[}8{]}}}Accusation and
  testimonies from the Old Testament (§3--16)

  \begin{enumerate}
  \tightlist
  \item
    \leavevmode\marginpar{\textcolor{parnum}{{[}9{]}}}Moses and the tent
    (§6--8)
  \item
    Solomon and the first temple (§9)
  \item
    Ezechiel and his vision (§10--13)
  \item
    Summary of passages from the Old Testament used for the defense of
    images (§14--16)
  \end{enumerate}
\item
  \leavevmode\marginpar{\textcolor{parnum}{{[}10{]}}}The New Testament
  (§17--20)

  \begin{enumerate}
  \tightlist
  \item
    \leavevmode\marginpar{\textcolor{parnum}{{[}11{]}}}Paul in Athens
    (§17--20)
  \end{enumerate}
\item
  \leavevmode\marginpar{\textcolor{parnum}{{[}12{]}}}Church fathers
  (§21--37)

  \begin{enumerate}
  \tightlist
  \item
    \leavevmode\marginpar{\textcolor{parnum}{{[}13{]}}}John of
    Constantinople (§21--23) (John Chrysostom)
  \item
    Image of Christ in Edessa (§24)
  \item
    Bishop Severianos (§25--26)
  \item
    Gregory the Illuminator (§27--29)
  \item
    Accusation of the opponents of images (§30--32)
  \item
    Eusebios' Church History (§33--37)
  \end{enumerate}
\item
  \leavevmode\marginpar{\textcolor{parnum}{{[}14{]}}}Accusations and
  comparisons (heathen Deities) (§38--43)

  \begin{enumerate}
  \tightlist
  \item
    \leavevmode\marginpar{\textcolor{parnum}{{[}15{]}}}Aramazd (§38--40)
  \item
    Anahit (§41--42)
  \item
    Astłik and Aphrodite (§43)
  \end{enumerate}
\item
  \leavevmode\marginpar{\textcolor{parnum}{{[}16{]}}}Other comparisons /
  materiality (§44--51)

  \begin{enumerate}
  \tightlist
  \item
    \leavevmode\marginpar{\textcolor{parnum}{{[}17{]}}}Pigment +
    writing/images (§44--45)
  \item
    The power of material objects (§46--48)
  \item
    Painted book-cases and their symbolic meaning (§49)
  \item
    Christ on a donkey as an example of representation (§50)
  \item
    Summary (§51)
  \end{enumerate}
\item
  \leavevmode\marginpar{\textcolor{parnum}{{[}18{]}}}Differences between
  Christian and heathen images (§52--66)

  \begin{enumerate}
  \tightlist
  \item
    \leavevmode\marginpar{\textcolor{parnum}{{[}19{]}}}Problem of
    representation (§57--58)
  \item
    The Armenian king Pap and the (Greek) tradition of images in
    churches (§59); concrete persons opposed to the use/veneration of
    images in churches
  \item
    Testimonies from the Old and New Testament (§60--61)
  \item
    Purity, pigment (matter), writing and images (§62--64)
  \item
    Paul and Timotheus/remembrance (§65--66)
  \end{enumerate}
\item
  \leavevmode\marginpar{\textcolor{parnum}{{[}20{]}}}End (quotes,
  accusation, prayer) (§67--70)
\end{enumerate}

\hypertarget{beispiel-fuxfcr-eine-tabelle}{%
\section{Beispiel für eine Tabelle}\label{beispiel-fuxfcr-eine-tabelle}}

\leavevmode\marginpar{\textcolor{parnum}{{[}21{]}}}Für die Aufstellung
siehe Tabelle \ref{tbl:tab1}.

\begin{longtable}[]{@{}llll@{}}
\caption{A test table \label{tbl:tab1}}\tabularnewline
\toprule
\begin{minipage}[b]{0.18\columnwidth}\raggedright
\textsc{Auspicious Statue/deities}\strut
\end{minipage} & \begin{minipage}[b]{0.13\columnwidth}\raggedright
\textsc{Moved from}\strut
\end{minipage} & \begin{minipage}[b]{0.13\columnwidth}\raggedright
\textsc{Moved to/resides in}\strut
\end{minipage} & \begin{minipage}[b]{0.44\columnwidth}\raggedright
\textsc{Sources/remarks}\strut
\end{minipage}\tabularnewline
\midrule
\endfirsthead
\toprule
\begin{minipage}[b]{0.18\columnwidth}\raggedright
\textsc{Auspicious Statue/deities}\strut
\end{minipage} & \begin{minipage}[b]{0.13\columnwidth}\raggedright
\textsc{Moved from}\strut
\end{minipage} & \begin{minipage}[b]{0.13\columnwidth}\raggedright
\textsc{Moved to/resides in}\strut
\end{minipage} & \begin{minipage}[b]{0.44\columnwidth}\raggedright
\textsc{Sources/remarks}\strut
\end{minipage}\tabularnewline
\midrule
\endhead
\begin{minipage}[t]{0.18\columnwidth}\raggedright
\emph{Śākyamuni Buddha}\strut
\end{minipage} & \begin{minipage}[t]{0.13\columnwidth}\raggedright
\strut
\end{minipage} & \begin{minipage}[t]{0.13\columnwidth}\raggedright
\strut
\end{minipage} & \begin{minipage}[t]{0.44\columnwidth}\raggedright
\strut
\end{minipage}\tabularnewline
\begin{minipage}[t]{0.18\columnwidth}\raggedright
釋迦牟尼佛

Śākyamuni Buddha\strut
\end{minipage} & \begin{minipage}[t]{0.13\columnwidth}\raggedright
靈鷲山

Gṛdhrakūṭa-parvata ``Vulture Peak''\strut
\end{minipage} & \begin{minipage}[t]{0.13\columnwidth}\raggedright
牛頭山

Oxhead Mountain\strut
\end{minipage} & \begin{minipage}[t]{0.44\columnwidth}\raggedright
S.2113v1-2: 釋迦牟尼佛從靈鷲山向牛頭山說法來 Śākyamuni Buddha came from
Gṛdhrakūṭa-parvata to Oxhead Mountain \emph{in order to proclaim the
Dharma}.

D-MG 231 (south): right:此牛頭山像從耆山履left:空而來

{[}Oxhead Mountain Tableaux in Dūnhuáng: D-MG 9, 340, 454; YL
33{]}\strut
\end{minipage}\tabularnewline
\begin{minipage}[t]{0.18\columnwidth}\raggedright
此像 (S.2114, P.3033)

``This statue''

憍賞彌國佛 (P.3352)

``The Buddha of the Country of Kauśāmbi''\strut
\end{minipage} & \begin{minipage}[t]{0.13\columnwidth}\raggedright
憍賞彌國

Country of Kauśāmbi\strut
\end{minipage} & \begin{minipage}[t]{0.13\columnwidth}\raggedright
于闐東媲摩城 (S.2114, P.3033)

Phema* east of Khotan

于闐 (P.3352)

Khotan\strut
\end{minipage} & \begin{minipage}[t]{0.44\columnwidth}\raggedright
S.2113v11-12: 此像從憍賞彌國飛往于闐東媲摩城今見在 Thus statue flew from
the country of Kauśāmbi to city of Phema in the east of Khotan
\emph{where it is present today}

P.3033v8-9: 此像從憍賞彌國飛往于闐東媲摩城中今見在殊靈瑞

P.3352r9: 憍賞彌國佛來往{[}=住{]}于闐

S.5659r9: 憍賞彌國佛來住于闐國

D-MG 231 (south): right: 于闐媲摩城中琱檀瑞像

D-MG 76\strut
\end{minipage}\tabularnewline
\begin{minipage}[t]{0.18\columnwidth}\raggedright
釋迦牟尼佛真容白檀身 The White sandal wood image created face-to-face
with (lit. ``true countenance'') Śākyamuni Buddha\strut
\end{minipage} & \begin{minipage}[t]{0.13\columnwidth}\raggedright
國王舍城

Rājagṛha\strut
\end{minipage} & \begin{minipage}[t]{0.13\columnwidth}\raggedright
于闐海眼寺

Ocean Eye Monastery of Khotan\strut
\end{minipage} & \begin{minipage}[t]{0.44\columnwidth}\raggedright
S.2113v16-17:
釋迦牟尼佛真容白檀身從國王舍城騰空而來在于闐海眼寺住{[}其像手把袈裟{]}
The white sandal statue of the true countenance of Śākyamuni Buddha was
soaring through the air from Rājagṛha and settled in the Ocean Eye
Monastery of Khotan {[}\emph{the hand of this statue grasps a kāṣāya}{]}

P.3353r2-3: 釋迦牟尼佛真容從國王舍城騰空而來在于闐海眼寺住

S.5659r5-6: 釋迦牟尼佛真容從國王舍城騰空而來在于闐海眼寺住

D-MG 72: 釋迦牟尼佛真容從國王舍城騰空如來在于闐海眼寺住

D-MG 45, 9, etc.

D-MG 231 (east): right:于闐海眼寺釋迦 left:聖容像 (``statue of sagely
countenance'')

D-MG 231 (north): right:釋迦牟尼真容從王舍 left:城騰空住海眼寺\strut
\end{minipage}\tabularnewline
\begin{minipage}[t]{0.18\columnwidth}\raggedright
釋迦牟尼佛真容白檀香為身

The true countenance of the Buddha Śākyamuni made of white sandalwood
(lit. white sandalwood being its body)\strut
\end{minipage} & \begin{minipage}[t]{0.13\columnwidth}\raggedright
漢國

Country of Hàn\strut
\end{minipage} & \begin{minipage}[t]{0.13\columnwidth}\raggedright
于闐坎城

Kǎnchéng of Khotan (= Phema?)\strut
\end{minipage} & \begin{minipage}[t]{0.44\columnwidth}\raggedright
S.2113v18-19:
釋迦牟尼佛真容白檀香**為身從漢國騰空而來在于闐坎城住{[}下其像手把袈裟{]}
The true countenance of the Buddha Śākyamuni made of white sandalwood
soared through the air form the Country of Hàn and came to settle at
Kănchéng {[}\emph{below, the statue's hand grasps a kāṣāya}{]}\strut
\end{minipage}\tabularnewline
\begin{minipage}[t]{0.18\columnwidth}\raggedright
釋迦牟尼佛

Śākyamuni Buddha\strut
\end{minipage} & \begin{minipage}[t]{0.13\columnwidth}\raggedright
舍衛國

Śrāvasti\strut
\end{minipage} & \begin{minipage}[t]{0.13\columnwidth}\raggedright
固城

``Old City'' {[}of Khotan{]}\strut
\end{minipage} & \begin{minipage}[t]{0.44\columnwidth}\raggedright
S.2113v20: 釋迦牟尼佛從舍衛國

騰空於固城住

{[}D-MG 231: right:于闐古城瑞像left:?{]}\strut
\end{minipage}\tabularnewline
\begin{minipage}[t]{0.18\columnwidth}\raggedright
釋迦牟尼

Śākyamuni\strut
\end{minipage} & \begin{minipage}[t]{0.13\columnwidth}\raggedright
舍衛國

Śrāvasti\strut
\end{minipage} & \begin{minipage}[t]{0.13\columnwidth}\raggedright
于闐國城

City of Khotan\strut
\end{minipage} & \begin{minipage}[t]{0.44\columnwidth}\raggedright
S.2113v31-32: 釋迦牟尼亦從舍衛國

騰空同來在于闐國城住{[}手把袈裟{]} Śākyamuni also soared through the air
from Śrāvasti and came to settle in City of Khotan {[}\emph{the hand
grasps a kāṣāya}{]}

D-MG 72 (caption nearly identically)\strut
\end{minipage}\tabularnewline
\begin{minipage}[t]{0.18\columnwidth}\raggedright
徽波陀佛

(not identified Buddha)\strut
\end{minipage} & \begin{minipage}[t]{0.13\columnwidth}\raggedright
舍衛國

Śrāvasti\strut
\end{minipage} & \begin{minipage}[t]{0.13\columnwidth}\raggedright
于闐城

City of Khotan\strut
\end{minipage} & \begin{minipage}[t]{0.44\columnwidth}\raggedright
S.2113v32-34: 徽波陀佛從舍衛國

住騰空而同來在于闐城住城人欽敬不可思議{[}其下像側{]} Huībōtuó Buddha
soared through the air from Śrāvasti and came together to settle in the
City of Khotan; the reverence of the inhabitants of the city was
inconceivable {[}\emph{below, the statue is slanted (?)}{]}

D-MG 76 (partly damaged; extant text identical)\strut
\end{minipage}\tabularnewline
\begin{minipage}[t]{0.18\columnwidth}\raggedright
\emph{Other buddhas}\strut
\end{minipage} & \begin{minipage}[t]{0.13\columnwidth}\raggedright
\strut
\end{minipage} & \begin{minipage}[t]{0.13\columnwidth}\raggedright
\strut
\end{minipage} & \begin{minipage}[t]{0.44\columnwidth}\raggedright
\strut
\end{minipage}\tabularnewline
\begin{minipage}[t]{0.18\columnwidth}\raggedright
迦葉佛

Kāśyapa Buddha (i.e., the Buddha prior to Śākyamuni)\strut
\end{minipage} & \begin{minipage}[t]{0.13\columnwidth}\raggedright
舍衛國

Śrāvasti\strut
\end{minipage} & \begin{minipage}[t]{0.13\columnwidth}\raggedright
于闐

Khotan\strut
\end{minipage} & \begin{minipage}[t]{0.44\columnwidth}\raggedright
S.2113v36-37:
迦葉佛亦從舍衛國騰空而來于闐國人皆虔敬不可思議{[}其像亦把袈裟{]} Kāśyapa
Buddha also soared through the air from Śrāvasti and came to Khotan; the
reverence {[}to the statue{]} of the people was inconceivable
{[}\emph{this statue also holds a kāṣāya}{]}

P.3353r7-8: 迦葉佛亦從舍衛國騰空而來在于闐國住國人虔敬無不遂願
{[}\ldots{}\emph{the inhabitants revered it and there was not one whose
wishes were not fulfilled}{]}

S.5659r10-11: 迦葉佛從舍衛國騰空而來在于闐國住國人虔敬無不遂願

D-MG 231: right: 迦葉佛從舍衛騰空left:於固城住\strut
\end{minipage}\tabularnewline
\begin{minipage}[t]{0.18\columnwidth}\raggedright
結迦宋佛

Kāśyapa Buddha (?)\strut
\end{minipage} & \begin{minipage}[t]{0.13\columnwidth}\raggedright
舍衛國

Śrāvasti\strut
\end{minipage} & \begin{minipage}[t]{0.13\columnwidth}\raggedright
固城

``Old City'' {[}of Khotan{]}\strut
\end{minipage} & \begin{minipage}[t]{0.44\columnwidth}\raggedright
S.2113v21: 結迦宋佛亦舍衛國來在固城住{[}其像手捻袈裟{]} Kāśyapa Buddha
(?) also came from Śrāvasti and settled in the Old City {[}\emph{the
hand of this statue holds the kāṣāya with its fingertips}{]}

D-MG 76 (caption partly damaged)

D-MG 231 (north)\strut
\end{minipage}\tabularnewline
\begin{minipage}[t]{0.18\columnwidth}\raggedright
Vipaśyin Buddha

毗波尸佛\strut
\end{minipage} & \begin{minipage}[t]{0.13\columnwidth}\raggedright
舍衛國

Śrāvasti\strut
\end{minipage} & \begin{minipage}[t]{0.13\columnwidth}\raggedright
于闐國

Country of Khotan\strut
\end{minipage} & \begin{minipage}[t]{0.44\columnwidth}\raggedright
P.3353r5: 毗婆尸佛{[}+從{]} 舍衛國

騰空而來在于闐國住有人欽仰不可思議 Vipaśyin Buddha soared through the
air from Śrāvasti and settled in the Country of Khotan; there were
people whose veneration was inconceivable

S.5659: 毗莎尸佛{[}+從{]} 舍衛國騰空而來在于闐國住有人欽仰不可思議

D-MG 231: (east)\strut
\end{minipage}\tabularnewline
\begin{minipage}[t]{0.18\columnwidth}\raggedright
伽儞迦牟尼佛

Jiānĭjiāmŏuní Buddha

(not identified Buddha)\strut
\end{minipage} & \begin{minipage}[t]{0.13\columnwidth}\raggedright
舍衛國

Śrāvasti\strut
\end{minipage} & \begin{minipage}[t]{0.13\columnwidth}\raggedright
固城

``Old City'' {[}of Khotan{]}\strut
\end{minipage} & \begin{minipage}[t]{0.44\columnwidth}\raggedright
S.2113v38-39: \sout{迦葉佛}
伽儞迦牟尼佛從舍衛國騰空而來在固城住{[}其像手捻袈裟{]} (note that there
is a strikethrough of the preceding迦葉佛Kāśyapa Buddha)

D-MG 76 (caption damaged)\strut
\end{minipage}\tabularnewline
\begin{minipage}[t]{0.18\columnwidth}\raggedright
\emph{Protector deities}\strut
\end{minipage} & \begin{minipage}[t]{0.13\columnwidth}\raggedright
\strut
\end{minipage} & \begin{minipage}[t]{0.13\columnwidth}\raggedright
\strut
\end{minipage} & \begin{minipage}[t]{0.44\columnwidth}\raggedright
\strut
\end{minipage}\tabularnewline
\begin{minipage}[t]{0.18\columnwidth}\raggedright
\strut
\end{minipage} & \begin{minipage}[t]{0.13\columnwidth}\raggedright
\strut
\end{minipage} & \begin{minipage}[t]{0.13\columnwidth}\raggedright
\strut
\end{minipage} & \begin{minipage}[t]{0.44\columnwidth}\raggedright
The standard phrase for all the protectors: ``Deity X protects the
country of Khotan.'' The Sanskrit transcriptions of the names of the
deities is based on Soymié 1984, 86. Protector deities and bodhisattvas
are depicted in D-MG 9, 98, 152, 154\strut
\end{minipage}\tabularnewline
\begin{minipage}[t]{0.18\columnwidth}\raggedright
迦迦那莎利神

Deity Gaganasvara\strut
\end{minipage} & \begin{minipage}[t]{0.13\columnwidth}\raggedright
\strut
\end{minipage} & \begin{minipage}[t]{0.13\columnwidth}\raggedright
\strut
\end{minipage} & \begin{minipage}[t]{0.44\columnwidth}\raggedright
S.2113v48-49: 迦迦那莎利神守護于闐國\strut
\end{minipage}\tabularnewline
\begin{minipage}[t]{0.18\columnwidth}\raggedright
莎那末利神

Deity Suvarṇamāla\strut
\end{minipage} & \begin{minipage}[t]{0.13\columnwidth}\raggedright
\strut
\end{minipage} & \begin{minipage}[t]{0.13\columnwidth}\raggedright
\strut
\end{minipage} & \begin{minipage}[t]{0.44\columnwidth}\raggedright
S.2113v49-50: 莎那末利神守護于闐國\strut
\end{minipage}\tabularnewline
\begin{minipage}[t]{0.18\columnwidth}\raggedright
莎那摩利神

Deity Suvarṇamāla\strut
\end{minipage} & \begin{minipage}[t]{0.13\columnwidth}\raggedright
\strut
\end{minipage} & \begin{minipage}[t]{0.13\columnwidth}\raggedright
\strut
\end{minipage} & \begin{minipage}[t]{0.44\columnwidth}\raggedright
S.2113v50: 莎那摩利神守護于闐國\strut
\end{minipage}\tabularnewline
\begin{minipage}[t]{0.18\columnwidth}\raggedright
阿隅闍天女

Heavenly Maiden Aṅkuśavatī\strut
\end{minipage} & \begin{minipage}[t]{0.13\columnwidth}\raggedright
\strut
\end{minipage} & \begin{minipage}[t]{0.13\columnwidth}\raggedright
\strut
\end{minipage} & \begin{minipage}[t]{0.44\columnwidth}\raggedright
S.2113v50: 阿隅闍天女守護闐國\strut
\end{minipage}\tabularnewline
\begin{minipage}[t]{0.18\columnwidth}\raggedright
毗沙天王神

Deity Vaiśravaṇa\strut
\end{minipage} & \begin{minipage}[t]{0.13\columnwidth}\raggedright
\strut
\end{minipage} & \begin{minipage}[t]{0.13\columnwidth}\raggedright
\strut
\end{minipage} & \begin{minipage}[t]{0.44\columnwidth}\raggedright
S.2113v51: 毗沙天王神守護于闐國\strut
\end{minipage}\tabularnewline
\begin{minipage}[t]{0.18\columnwidth}\raggedright
阿婆羅質多神

Deity Apārajita\strut
\end{minipage} & \begin{minipage}[t]{0.13\columnwidth}\raggedright
\strut
\end{minipage} & \begin{minipage}[t]{0.13\columnwidth}\raggedright
\strut
\end{minipage} & \begin{minipage}[t]{0.44\columnwidth}\raggedright
S.2113v51: 阿婆羅質多神守護于闐國

P.3352r2: 阿婆羅質多神護于闐國

S.5659r4: 阿婆羅質多神護于闐國\strut
\end{minipage}\tabularnewline
\begin{minipage}[t]{0.18\columnwidth}\raggedright
摩訶迦羅神

Deity Mahākāla\strut
\end{minipage} & \begin{minipage}[t]{0.13\columnwidth}\raggedright
\strut
\end{minipage} & \begin{minipage}[t]{0.13\columnwidth}\raggedright
\strut
\end{minipage} & \begin{minipage}[t]{0.44\columnwidth}\raggedright
S.2113v52: 摩訶迦羅神守護于闐國\strut
\end{minipage}\tabularnewline
\begin{minipage}[t]{0.18\columnwidth}\raggedright
悉他那天女

恭陀那天女

Heavenly Maiden Sthānavatī\strut
\end{minipage} & \begin{minipage}[t]{0.13\columnwidth}\raggedright
\strut
\end{minipage} & \begin{minipage}[t]{0.13\columnwidth}\raggedright
\strut
\end{minipage} & \begin{minipage}[t]{0.44\columnwidth}\raggedright
S.2113v52: 悉他那天女護于闐國

S.5659r8: 恭陀那天女守護于闐國\strut
\end{minipage}\tabularnewline
\begin{minipage}[t]{0.18\columnwidth}\raggedright
\emph{Protector Bodhisattvas}\strut
\end{minipage} & \begin{minipage}[t]{0.13\columnwidth}\raggedright
\strut
\end{minipage} & \begin{minipage}[t]{0.13\columnwidth}\raggedright
\strut
\end{minipage} & \begin{minipage}[t]{0.44\columnwidth}\raggedright
\strut
\end{minipage}\tabularnewline
\begin{minipage}[t]{0.18\columnwidth}\raggedright
\strut
\end{minipage} & \begin{minipage}[t]{0.13\columnwidth}\raggedright
\strut
\end{minipage} & \begin{minipage}[t]{0.13\columnwidth}\raggedright
\strut
\end{minipage} & \begin{minipage}[t]{0.44\columnwidth}\raggedright
Same formula as in \emph{Protector deities}.\strut
\end{minipage}\tabularnewline
\begin{minipage}[t]{0.18\columnwidth}\raggedright
金剛藏菩薩

Vajragarbha Bodhisattva\strut
\end{minipage} & \begin{minipage}[t]{0.13\columnwidth}\raggedright
\strut
\end{minipage} & \begin{minipage}[t]{0.13\columnwidth}\raggedright
\strut
\end{minipage} & \begin{minipage}[t]{0.44\columnwidth}\raggedright
S.5659r1: 金剛藏菩薩護于闐國\strut
\end{minipage}\tabularnewline
\begin{minipage}[t]{0.18\columnwidth}\raggedright
觀世音菩薩

Avalokiteśvara Bodhisattva\strut
\end{minipage} & \begin{minipage}[t]{0.13\columnwidth}\raggedright
\strut
\end{minipage} & \begin{minipage}[t]{0.13\columnwidth}\raggedright
\strut
\end{minipage} & \begin{minipage}[t]{0.44\columnwidth}\raggedright
S.5659r14: 觀世音菩薩助于闐國\strut
\end{minipage}\tabularnewline
\begin{minipage}[t]{0.18\columnwidth}\raggedright
寶壇花菩薩

Bǎotánhuā púsà\strut
\end{minipage} & \begin{minipage}[t]{0.13\columnwidth}\raggedright
\strut
\end{minipage} & \begin{minipage}[t]{0.13\columnwidth}\raggedright
\strut
\end{minipage} & \begin{minipage}[t]{0.44\columnwidth}\raggedright
S.5659r14-15: 寶壇花菩薩助于闐國

(No reference to this Bodhisattva other than in a caption of D-MG
449)\strut
\end{minipage}\tabularnewline
\begin{minipage}[t]{0.18\columnwidth}\raggedright
\emph{``Narratives'' mentioning Khotan}\strut
\end{minipage} & \begin{minipage}[t]{0.13\columnwidth}\raggedright
\strut
\end{minipage} & \begin{minipage}[t]{0.13\columnwidth}\raggedright
\strut
\end{minipage} & \begin{minipage}[t]{0.44\columnwidth}\raggedright
\strut
\end{minipage}\tabularnewline
\begin{minipage}[t]{0.18\columnwidth}\raggedright
于闐王河浴佛瑞像

The King of Khotan bathing Buddha's Auspicious Statue in the river\strut
\end{minipage} & \begin{minipage}[t]{0.13\columnwidth}\raggedright
\strut
\end{minipage} & \begin{minipage}[t]{0.13\columnwidth}\raggedright
\strut
\end{minipage} & \begin{minipage}[t]{0.44\columnwidth}\raggedright
S.2113v22-23: 于闐王河浴佛瑞像身丈餘杖錫持鉢畫形布(?)立{[}其像赤體立{]}
The King of Khotan bathes the Auspicious Statue of the Buddha in the
river; the body {[}of the statue{]} is more than {[}one{]} \emph{zhàng}
{[}in height; i.e.~more than 3 meters{]}; he holds a monk's staff and
begging bowl; \ldots{} (?) {[}\emph{this statue is naked and stands
upright}{]}\strut
\end{minipage}\tabularnewline
\begin{minipage}[t]{0.18\columnwidth}\raggedright
于闐國舍利弗毗沙門天王

Śāriputra and the Heavenly King Vaiśravaṇa of Khotan\strut
\end{minipage} & \begin{minipage}[t]{0.13\columnwidth}\raggedright
\strut
\end{minipage} & \begin{minipage}[t]{0.13\columnwidth}\raggedright
\strut
\end{minipage} & \begin{minipage}[t]{0.44\columnwidth}\raggedright
P.3352r11-12: 于闐國舍利弗毗沙門天王決海時

The scene when Śāriputra and the Heavenly King Vaiśravaṇa of Khotan
divide the ocean

D-MG 9, 237, 454

D-MG 231: right:于闐國舍利弗毗沙門天王決海時\strut
\end{minipage}\tabularnewline
\bottomrule
\end{longtable}

\hypertarget{beispiel-fuxfcr-ein-bild}{%
\section{Beispiel für ein Bild}\label{beispiel-fuxfcr-ein-bild}}

\leavevmode\marginpar{\textcolor{parnum}{{[}22{]}}}Der Ausstellungsraum
ist in Abbildung \ref{fig:img1} zu sehen.

\begin{figure}
\hypertarget{fig:img1}{%
\centering
\includegraphics{Image 1.tif}
\caption{Insert Image 1 here}\label{fig:img1}
}
\end{figure}

\hypertarget{bibliography}{%
\section*{References}\label{bibliography}}
\addcontentsline{toc}{section}{References}

\hypertarget{refs}{}
\begin{cslreferences}
\leavevmode\hypertarget{ref-beagle_digital_2014}{}%
Beagle, Donald. 2014. ``Digital Humanities in the Research Commons:
Precedents \& Prospects.'' Text. \emph{Dh+lib} (blog). January 30, 2014.
\url{https://acrl.ala.org/dh/2014/01/30/digital-humanities-in-the-research-commons-precedents-prospects-3/}.

\leavevmode\hypertarget{ref-berry_understanding_2012}{}%
Berry, David M. 2012. \emph{Understanding Digital Humanities}.
Basingstoke: Palgrave Macmillan.
\url{https://www.palgrave.com/de/book/9780230292642}.

\leavevmode\hypertarget{ref-blume_burkhard_2016}{}%
Blume, Michael. 2016. ``Burkhard Gladigow und die europäische
Religionsgeschichte von der griechischen Antike bis zu den Protokollen
der Weisen von Zion.'' \emph{Natur des Glaubens} (blog). June 17, 2016.
\url{https://scilogs.spektrum.de/natur-des-glaubens/burkhard-gladigow-religionsgeschichte-antike-protokollen/}.

\leavevmode\hypertarget{ref-burdick_digital_humanities_2012}{}%
Burdick, Anne, Johanna Drucker, Peter Lunenfeld, Todd Presner, and
Jeffrey Schnapp. 2012. \emph{Digital\_Humanities}. MIT Press.

\leavevmode\hypertarget{ref-busa_annals_1980}{}%
Busa, R. 1980. ``The Annals of Humanities Computing: The Index
Thomisticus.'' \emph{Computers and the Humanities} 14 (2): 83--90.
\url{https://doi.org/10.1007/BF02403798}.

\leavevmode\hypertarget{ref-celentano_informatica_2004}{}%
Celentano, Augusto, Agostino Cortesi, and Paolo Mastandrea. 2004.
``Informatica Umanistica: Una Disciplina Di Confine.'' \emph{Mondo
Digitale} 4: 44--55.

\leavevmode\hypertarget{ref-classen_towards_2012}{}%
Classen, Christoph, Susanne Kinnebrock, and Maria Löblich. 2012.
``Towards Web History: Sources, Methods, and Challenges in the Digital
Age. An Introduction.'' \emph{Historical Social Research / Historische
Sozialforschung} 37 (4 (142)): 97--101.
\url{https://www.jstor.org/stable/41756476}.

\leavevmode\hypertarget{ref-gladigow_europaische_1995}{}%
Gladigow, Burkhard. 1995. ``Europäische Religionsgeschichte.'' In
\emph{Lokale Religionsgeschichte}, edited by G. Hans Kippenberg and
Brigitte Luchesi, 21--42. Marburg: Diagonal.

\leavevmode\hypertarget{ref-isomae_religious_2014}{}%
Isomae, Jun'ichi. 2014. \emph{Religious Discourse in Modern Japan:
Religion, State, and Shintō}. Translated by Galen Amstutz and Lynne E.
Riggs. London: Brill.

\leavevmode\hypertarget{ref-kippenberg_europaische_2009}{}%
Kippenberg, Hans G., Jörg Rüpke, and Kocku von Stuckrad, eds. 2009.
\emph{Europäische Religionsgeschichte: Ein Mehrfacher Pluralismus}.
Göttingen: Vandenhoeck \& Ruprecht.

\leavevmode\hypertarget{ref-rudolph_leipzig_1962}{}%
Rudolph, Kurt. 1962. ``Leipzig Und Die Religionswissenschaft.''
\emph{Numen} 9 (1): 53--68. \url{https://doi.org/10.2307/3269414}.

\leavevmode\hypertarget{ref-schalk_religion_2013}{}%
Schalk, Peter, ed. 2013. \emph{Religion in Asien? Studien Zur
Anwendbarkeit Des Religionsbegriffs}. Acta Universitatis Upsaliensis,
Historia Religionum 32. Uppsala: Uppsala Universitet.

\leavevmode\hypertarget{ref-seiwert_popular_2003}{}%
Seiwert, Hubert. 2003. \emph{Popular Religious Movements and Heterodox
Sects in Chinese History}. Leiden/Boston: Brill.

\leavevmode\hypertarget{ref-waardenburg_islam:_2002}{}%
Waardenburg, Jean Jacques. 2002. \emph{Islam: Historical, Social, and
Political Perspectives}. Religion and Reason 40. Berlin: de Gruyter.

\leavevmode\hypertarget{ref-waardenburg_muslims_2007}{}%
---------. 2007. \emph{Muslims as Actors: Islamic Meanings and Muslim
Interpretations in the Perspective of the Study of Religions}. Religion
and Reason 46. Berlin: de Gruyter.

\leavevmode\hypertarget{ref-weber_gesammelte_1988}{}%
Weber, Max. (1922) 1988. \emph{Gesammelte Aufsätze Zur
Religionssoziologie}. 9th ed. Vol. 2. Tübingen: Mohr.
\end{cslreferences}

\end{document}
